\chapter{Boot Process}
\label{ch:boot}

This chapter traces the execution path from power-on to the login prompt,
detailing every step the system takes from firmware initialization through
kernel boot to userspace startup.

\section{Boot Sequence Overview}

The complete boot sequence has five major phases:

\begin{enumerate}[nosep]
  \item \textbf{Firmware}: SeaBIOS POST, device enumeration, boot device
        selection.
  \item \textbf{Stage 1 Bootloader}: GRUB boot.img loads core.img from
        the MBR gap sectors.
  \item \textbf{Stage 2 Bootloader}: GRUB core.img loads the ext2 driver,
        reads grub.cfg, and loads the kernel.
  \item \textbf{Kernel}: Decompression, hardware initialization, initramfs
        extraction, PID 1 execution.
  \item \textbf{Userspace}: BusyBox init runs rcS, which mounts filesystems,
        configures networking, and starts services.
\end{enumerate}

\section{Phase 1: Firmware (SeaBIOS)}

\subsection{Power-On Self Test}

When the virtual machine starts, SeaBIOS (shipped with QEMU) executes:

\begin{enumerate}[nosep]
  \item CPU reset vector: The CPU starts executing at address 0xFFFFFFF0
        (the reset vector) in real mode. SeaBIOS jumps to its entry point
        at 0xE05B.
  \item Self-test: The BIOS checks the CPU, memory controller, and DMA
        controller. On QEMU, these checks are minimal since the hardware
        is emulated.
  \item Memory sizing: The BIOS detects the RAM size by probing the
        memory controller. QEMU reports the memory size configured via
        the \texttt{-m} option.
  \item PCI enumeration: The BIOS scans the PCI bus (configuration space
        accessed via I/O ports 0xCF8/0xCFC) to discover VirtIO devices,
        the VGA adapter, and other peripherals.
  \item ACPI table generation: The BIOS generates the RSDT, DSDT, FACP,
        APIC, and HPET tables describing the virtual hardware
        configuration.
  \item Boot device selection: The BIOS iterates through the boot order
        (set by QEMU's \texttt{-boot} option). For each device, it reads
        LBA 0 into memory at address 0x7C00. If the last two bytes are
        0x55 0xAA, it transfers control.
\end{enumerate}

\section{Phase 2: GRUB Stage 1}

At address 0x7C00, GRUB's \texttt{boot.img} begins execution:

\begin{enumerate}[nosep]
  \item The CPU is in real mode with the following register state:
        \texttt{DL = boot drive number}, \texttt{CS:IP = 0x0000:0x7C00}.
  \item \texttt{boot.img} sets up segment registers (\texttt{DS, ES, SS})
        to 0x0000 and initializes the stack at 0x7C00 (growing downward).
  \item It reads the blocklist from \texttt{core.img} at LBA 1. The
        blocklist (12 bytes at offset 0x1F4) specifies how many sectors
        to load and where.
  \item Using INT 13h function 42h (extended read), it loads
        \texttt{diskboot.img} (the first sector of core.img) into memory
        at the segment specified by the blocklist.
  \item Execution jumps to the loaded diskboot sector.
\end{enumerate}

\section{Phase 3: GRUB Stage 2}

\subsection{diskboot.img}

The diskboot sector performs:

\begin{enumerate}[nosep]
  \item Switches CPU from real mode to protected mode by setting bit 0
        of CR0.
  \item Enables the A20 gate (via port 0x92 or the keyboard controller
        method) to allow access to memory above 1 MB.
  \item Sets up a minimal Global Descriptor Table (GDT) with flat
        segmentation (base 0, limit 4 GB).
  \item Loads the rest of core.img using the blocklist.
  \item Jumps to the GRUB kernel entry point.
\end{enumerate}

\subsection{GRUB Kernel Initialization}

The GRUB kernel initializes its internal subsystems:

\begin{enumerate}[nosep]
  \item Memory management: GRUB detects the system memory map (using
        INT 15h function E820h) and sets up its internal allocator.
  \item Device manager: Probes disks via BIOS INT 13h and partitions
        via the part\_msdos module.
  \item Filesystem driver: The embedded ext2 module reads the VEILBOX
        partition's superblock and root directory.
  \item Config file parsing: GRUB reads \texttt{/boot/grub/grub.cfg},
        displays the menu (with 5-second timeout), and executes the
        selected menuentry.
  \item Kernel loading: The \texttt{linux} command loads the kernel
        (\texttt{/boot/vmlinuz}) into memory using the 32-bit boot
        protocol. It sets up the kernel command line and jumps to the
        kernel's 32-bit entry point.
\end{enumerate}

\section{Phase 4: Kernel Initialization}

\subsection{Decompression}

The kernel begins at \texttt{arch/x86/boot/compressed/head\_64.S} which:

\begin{enumerate}[nosep]
  \item Sets up the 64-bit long mode page tables (identity mapping).
  \item Enables PAE (Physical Address Extension) and long mode by
        setting EFER.LME in the MSR.
  \item Jumps to the 64-bit C code (\texttt{extract\_kernel()}).
  \item Decompresses the kernel proper (from the gzip or xz-compressed
        data in the .init section) to the target address.
  \item Jumps to the decompressed kernel entry point
        (\texttt{start\_kernel()}).
\end{enumerate}

\subsection{start\_kernel()}

The architecture-independent kernel entry:

\begin{enumerate}[nosep]
  \item Architecture setup: \texttt{setup\_arch()} initializes the page
        table, IDT, GDT, TSS, and APIC.
  - Memory management: \texttt{mm\_init()} sets up the buddy allocator,
        slab allocator, and vmalloc region.
  - Scheduler initialization: \texttt{sched\_init()} creates the idle
        thread and initializes runqueues.
  - Interrupt subsystem: \texttt{init\_IRQ()} programs the APIC and
        installs interrupt handlers.
  - Timers: \texttt{time\_init()} sets up the HPET or TSC as the
        clocksource.
  - Console: \texttt{console\_init()} registers the serial console
        driver (8250 UART) and VGA text console.
  - Initramfs: \texttt{populate\_rootfs()} extracts the embedded cpio
        archive to tmpfs.
  - PID 1: \texttt{run\_init\_process()} locates /sbin/init and executes
        it as PID 1.
\end{enumerate}

\section{Phase 5: Userspace Initialization}

BusyBox init (\texttt{/sbin/init}) runs as PID 1:

\begin{enumerate}[nosep]
  \item Reads \texttt{/etc/inittab} to determine what actions to take.
  \item Executes the \texttt{::sysinit:/etc/init.d/rcS} action, running
        the system initialization script.
  \item rcS mounts proc, sysfs, devtmpfs, devpts, and tmpfs.
  \item rcS sets the hostname and mounts the state partition.
  \item rcS configures the loopback interface and starts DHCP on eth0.
  \item rcS starts syslogd, containerd, and Dropbear.
  \item After rcS completes, init spawns getty on tty1 and ttyS0.
  \item Login prompt appears: \texttt{veilbox login:}.
\end{enumerate}


\section{Kernel Command-Line Parameters}

The kernel command line is passed via the GRUB configuration or QEMU's
\texttt{-append} option. Key parameters used by Veilbox:

\begin{longtable}{|p{4cm}|p{3cm}|p{5cm}|}
\hline
\textbf{Parameter} & \textbf{Value} & \textbf{Purpose} \\
\hline
console & ttyS0,tty1 & Set serial and VGA console \\
root & /dev/ram0 & Root device (initramfs) \\
rw & (implied) & Mount root read-write \\
quiet & (optional) & Reduce boot messages \\
veilbox.autologin & (flag) & Enable auto-login \\
panic & 10 & Reboot after 10s on panic \\
init & /sbin/init & Path to init process \\
\hline
\end{longtable}

\subsection{Debug Parameters}

For troubleshooting boot issues:

\begin{lstlisting}
# Maximum verbosity
console=ttyS0 earlyprintk=serial,ttyS0,115200 debug ignore_loglevel

# Break before mounting root
rd.break

# Single-user mode
single

# Skip init entirely (emergency shell)
init=/bin/sh
\end{lstlisting}

\subsection{The veilbox.autologin Parameter}

This custom parameter is parsed by the \texttt{/sbin/autologin} wrapper.
It checks \texttt{/proc/cmdline} for the string:

\begin{lstlisting}[language=bash]
#!/bin/sh
if grep -q veilbox.autologin /proc/cmdline; then
    # Auto-login as root without password
    login -f root
else
    # Normal getty with password prompt
    exec /sbin/getty "$@"
fi
\end{lstlisting}

\section{Sysfs and Procfs Initialization}

After mounting proc and sysfs, the kernel exposes runtime information:

\begin{lstlisting}
# Process information
/proc/cpuinfo     # CPU details (model, flags, cores)
/proc/meminfo     # Memory usage details
/proc/self/ns     # Namespace IDs for current process
/proc/1/root      # Root filesystem of PID 1

# Kernel parameters
/proc/sys/kernel/  # Kernel tuning parameters
/proc/sys/net/     # Network stack tuning

# Device information (sysfs)
/sys/class/block/  # Block devices
/sys/class/net/    # Network interfaces
/sys/devices/      # Device hierarchy
\end{lstlisting}
